%! Justifying a Claim About the Difference of Two Means Based on a CI — Unit 7, Lesson 7.7 — Warm-Up
\documentclass[10pt]{article}
\usepackage{apstats-article}
\usepackage{apstats-boxes}
\newcommand{\TallMath}[1]{$\displaystyle #1\rule[-1.4em]{0pt}{3.2em}$}

\begin{document}

\pageheader{Unit 7, Lesson 7.7}{Warm-Up}
\namedateperiod

\vspace{0.15in}

\begin{notesbox}{Spiral Review --- Two-Sample $t$-Interval \& CI Interpretation}

Two researchers study daily screen time (hours) for teenagers in two regions.
Region 1: $\bar x_1 = 6.4$ hr, $s_1 = 1.8$ hr, $n_1 = 30$.
Region 2: $\bar x_2 = 5.7$ hr, $s_2 = 2.1$ hr, $n_2 = 35$.
Both conditions for inference are met. $df = 58$ (given). $t^* = 2.002$ for 95\% confidence.

\medskip
\textbf{1.} Compute the standard error $SE = \sqrt{\dfrac{s_1^2}{n_1} + \dfrac{s_2^2}{n_2}}$.

\[SE = \sqrt{\dfrac{\blank{1.8cm}^2}{\blank{1.5cm}} + \dfrac{\blank{1.8cm}^2}{\blank{1.5cm}}} = \sqrt{\blank{2.5cm}} \approx \blank{2cm}\]

\medskip
\textbf{2.} Construct the 95\% CI for $\mu_1 - \mu_2$.

\[(\bar x_1 - \bar x_2) \pm t^* \cdot SE = (\blank{1.5cm} - \blank{1.5cm}) \pm \blank{2cm} \cdot \blank{2cm} \approx (\blank{2cm},\; \blank{2cm})\text{ hr}\]

\medskip
\textbf{3.} Interpret the interval in context.

\writelines{2}

\medskip
\textbf{4.} Based on the interval, is there convincing evidence that teenagers in Region 1
have a higher mean screen time than those in Region 2? Circle one and explain briefly.

\quad \textbf{Yes} \qquad \textbf{No} \qquad \textbf{Cannot tell}

\writeline

\end{notesbox}

\end{document}
